% Accented characters: \usepackage[utf8]{inputenc} is no longer needed.
% UTF-8 has been the default input encoding of LaTeX since 2018, and
% LuaLaTeX reads UTF-8 natively.

% CJK (Chinese, Japanese, Korean) characters: \cjk{改善}
% CJKutf8 only works with pdfLaTeX, so under LuaLaTeX (the engine we use
% for tagged, accessible PDFs) CJK text is set in the Harano Aji Mincho
% font that ships with TeX Live (and therefore with Overleaf).
% Final mode (TAPS) loads neither iftex nor fontspec (acmart loads both
% under LuaLaTeX) nor CJKutf8 (not on ACM's list), so there \cjk needs
% LuaLaTeX.
\ifchizenfinal
\ifLuaTeX
  \newfontfamily\chizenCJKfont{HaranoAjiMincho-Regular.otf}
  \newcommand{\cjk}[1]{{\chizenCJKfont #1}}
\else
  \newcommand{\cjk}[1]{\PackageError{CHIZen}{\string\cjk\space needs LuaLaTeX in final mode}{%
    CJKutf8 is not on ACM's TAPS list. Compile with LuaLaTeX.}}
\fi
\else
\usepackage{iftex}
\ifLuaTeX
  \usepackage{fontspec}
  \newfontfamily\chizenCJKfont{HaranoAjiMincho-Regular.otf}
  \newcommand{\cjk}[1]{{\chizenCJKfont #1}}
\else
  \usepackage{CJKutf8}
  \newcommand{\cjk}[1]{\begin{CJK}{UTF8}{min}#1\end{CJK}}
\fi
\fi


% ACCESSIBILITY
% The old `accessibility' package is obsolete and must not be loaded.

% MARKDOWN (draft mode only: ACM's TAPS does not accept it)
%
% TODO: Explore how useful this is for general-purpose HCI papers
%       - Tables/Figures need to be a separate thing, but useful for everything else?
%
% markdown v3 loads amssymb, which clashes with the newtxmath fonts that
% acmart already loads (both define \Bbbk etc.). newtxmath provides all the
% amssymb symbols, so tell LaTeX amssymb is already there.
\ifchizenfinal\else
\makeatletter
\@ifpackageloaded{newtxmath}{\@namedef{ver@amssymb.sty}{2013/01/14 v3.01 (symbols provided by newtxmath)}}{}
\makeatother
% `notes' was called `footnotes' before markdown v3.
% LaTeX inside a markdown block: maths between dollar signs (texMathDollars),
% any other command as a raw span, `\citep{key}`{=tex} (rawAttribute), and
% % comments are dropped (texComments). These replace the `hybrid' option,
% which markdown 3.7.1 soft-deprecated. Plain text is now escaped, so write
% & as it is (not \\&); a % still starts a comment, so write \% for a percent.
\usepackage[notes,definitionLists,hashEnumerators,smartEllipses,rawAttribute,texMathDollars,texComments]{markdown}

\usepackage{paracol}
\usepackage{minted}
% fvextra (under minted) moves each code line up by 0.25pt with \vspace, so
% that background colours of adjacent lines overlap; acmart warns about any
% \vspace in the text. Listings here have no background colour, so turn it
% off. (With bgcolor=..., add bgcolorboxoverlap=0.25pt to that listing.)
\fvset{backgroundcolorboxoverlap=0pt}
% acmart's monospaced font, Inconsolata, has no italic, which minted styles
% use for comments. Say so (set them upright, as LaTeX would after a "Font
% shape undefined" warning). No spaces in the size specification: it is
% read from a macro.
\makeatletter
\begingroup\ttfamily\endgroup % declares the font family
\ifcsname\encodingdefault/\ttdefault/m/it\endcsname\else
  \edef\chizen@tmp{\noexpand\DeclareFontShape{\encodingdefault}{\ttdefault}{m}{it}{<->ssub*\ttdefault/m/n}{}}%
  \chizen@tmp
\fi
\makeatother
\fi
\ifchizenfinal
% Final mode: markdown, paracol and minted are not on ACM's TAPS list.
% Using them stops with an error; the body of the environment is skipped
% (with the comment package, which acmart loads, so verbatim text is fine).
\makeatletter
\def\chizen@textnotinfinal#1#2{%
  \PackageError{CHIZen}{#1 is not available in final mode}{%
    ACM's TAPS does not accept #1, so CHI Zen does not load it when the
    paper says \string\def\string\chizenmode{final}. #2 See the CHI Zen
    manual, Section "Camera-Ready for ACM TAPS".}}
\excludecomment{markdown}
\AddToHook{env/markdown/begin}{\chizen@textnotinfinal{markdown}{Rewrite
  the Markdown block in LaTeX.}}
\NewDocumentCommand\markdownInput{O{}m}{\chizen@textnotinfinal{markdown}{Rewrite
  #2 in LaTeX.}}
\excludecomment{minted}
\AddToHook{env/minted/begin}{\chizen@textnotinfinal{minted}{Use the
  listings package (\string\usepackage{listings} and a lstlisting
  environment) instead.}}
\NewDocumentCommand\mint{}{\chizen@textnotinfinal{minted (\string\mint)}{Use
  \string\lstinline\space from the listings package instead.}}
\NewDocumentCommand\mintinline{}{\chizen@textnotinfinal{minted (\string\mintinline)}{Use
  \string\lstinline\space from the listings package instead.}}
\NewDocumentCommand\inputminted{O{}mm}{\chizen@textnotinfinal{minted (\string\inputminted)}{Use
  \string\lstinputlisting\space from the listings package instead.}}
\NewDocumentCommand\usemintedstyle{O{}m}{}
\NewDocumentEnvironment{paracol}{m +b}{\chizen@textnotinfinal{paracol}{Use
  a table or acmart's two-column layout instead.}}{}
\makeatother
\fi
\usepackage{fontawesome5}

% This package is useful for revisions! 
% It defines \comment, which acmart's comment package defines too: keep
% comment's \comment and rename changes' one to \SCcomment.
\ifchizenfinal
% Final mode: savesym is not on ACM's TAPS list, so do the same by hand,
% and accept every change ([final]: no colours or struck-out text in the
% camera-ready version). commentmarkup=footnote keeps changes from loading
% todonotes, which loads TikZ.
\let\chizenSavedComment\comment
\let\comment\undefined
\usepackage[final,commentmarkup=footnote]{changes}
\let\SCcomment\comment
\let\comment\chizenSavedComment
\else
\usepackage{savesym}
\savesymbol{comment}
\let\comment\undefined % savesym leaves \comment as \relax, which changes v4 treats as defined
\usepackage{changes}
%\usepackage[final]{changes}
\restoresymbol{SC}{comment}
\fi


% Mark likes to add these options to his papers .... worth considering
\frenchspacing
\raggedbottom

% Lennart's acronym package
\usepackage[nolist]{acronym}

\begin{acronym}
\acro{ANOVA}{Analysis of Variance}
\acro{HCI}{Human-Computer Interaction}
\end{acronym}


% Jim's quotation stuff
%\renewenvironment{quote}{%
%  \list{}{%
%    \leftmargin0.4cm   % this is the adjusting screw
%    \rightmargin\leftmargin
%  }
%  \item\relax
%}
%{\endlist}

% Jim added this -- useful highlighting function \hl{}
%\usepackage{soul,color}



% Study Variables
% ----------------------------

% scale up the sc font slightly
\ifchizenfinal
% Final mode: letltxmacro, lipsum and relsize are not on ACM's TAPS list.
% \NewCommandCopy is built into LaTeX; \chizenrelscale does what
% \relscale does.
\makeatletter
\NewCommandCopy{\oldscshape}{\scshape}
\def\chizenrelscale#1{%
  \@tempdima=#1\dimexpr\f@size pt\relax
  \@tempdimb=#1\dimexpr\f@baselineskip\relax
  \edef\chizen@tmp{\noexpand\fontsize{\the\@tempdima}{\the\@tempdimb}}%
  \chizen@tmp\selectfont}
\renewcommand{\scshape}[1]{\oldscshape\chizenrelscale{1.2}#1}
\NewDocumentCommand\lipsum{O{}}{\PackageError{CHIZen}{\string\lipsum\space
  is not available in final mode}{Replace the placeholder text with your own.}}
\makeatother
\else
\usepackage{letltxmacro,lipsum,relsize}
\LetLtxMacro{\oldscshape}{\scshape}
\renewcommand{\scshape}[1]{\oldscshape\relscale{1.2}#1}
\fi

% an experiment factor/variable
\newcommand{\f}[1]{\mbox{{\textsc{#1}}}\xspace}
% usage: \f{block}, \f{width}, etc.
% NOTE: need \xspace package to fix spacing issues
%       \mbox prevents breaking

% a code when qualitative coding 
\newcommand{\cd}[1]{\textit{``{#1}''}\xspace}
% usage: \cd{errors}, \cd{work}